\documentclass[border = 0mm]{standalone}

\usepackage[utf8]{inputenc}
\usepackage{fontspec}
\usepackage[default]{fontsetup}
\usepackage{amsmath}

\usepackage{xcolor}
% --
\usepackage{tikz}
\usetikzlibrary{calc,math,arrows.meta,angles,scopes,backgrounds}

\usepackage{tikz-3dplot}
\usetikzlibrary{3d}

\newcommand{\pointer}[1]{\node[red] at #1 {$\times$};}

% ---
\begin{document}

\begin{tikzpicture}[
	scale = 0.85,
	line join = round, line cap = round,
]

	\tikzmath{
		\Width = 5.6cm; \Height = 3/4*\Width;
	};

	\node[
		inner sep = 0mm,
		minimum width = \Width,
		minimum height = \Height,	
	] (N) at (0,0) {};
	
	\clip (N.north west) rectangle (N.south east);

	\begin{scope}[
		shift = (N),
		shift = {(0,-0.75cm)},
	]
	
	% set point of view
	\tdplotsetmaincoords{70}{120}
	
	% apply point of view
	\begin{scope}[
		tdplot_main_coords
	]
	
	% -- shift
	{ [shift = {(0cm, -0.2cm)}]
	
	% origin
	\coordinate (O) at (0,0,0);
	
	% axes
%	{ [gray!50]
%  	\draw (O) -- (3,0,0) node[black, shift = {(-1.4mm,-1.4mm)}] {$x$};
%  	\draw (O) -- (0,3,0) node[black, shift = {(2mm,-0.8mm)}] {$y$};
%  	\draw (O) -- (0,0,3) node[black, shift = {(0mm,2mm)}] {$z$};
%	}	
	
	% frustum dimensions
	\tikzmath{
		\R = 2.3; \r = 1; \H = 2.4;
		\A = \R*cos(45); \a = \r*cos(45);
	}
	
	% define points of a trucated pyramid
	\begin{scope}[canvas is xy plane at z = 0]
	
	\foreach \i in {1,...,4}
		\coordinate (P\i) at ({45 + 90*(\i - 1)}:\R); 
		
	\end{scope}
	
	\begin{scope}[canvas is xy plane at z = \H]
	\foreach \i in {1,...,4}
		\coordinate (p\i) at ({45 + 90*(\i - 1)}:\r); 
	\end{scope}
	
	% draw
	\draw[gray!50]
		(P3) edge (p3) edge (P4) edge (P2);
	
	\draw
		(p1)--(p2)--(p3)--(p4)--cycle
		(P4) edge (p4) -- (P1) edge (p1) -- (P2) edge (p2);	
	
	% shadows
	{ [on background layer]
	
	\fill[gray, opacity = 0.1]
		(p1)--(p2)--(p3)--(p4)
		(p4)--(P4)--(P1)--(p1)--cycle;
		
	\fill[gray, opacity = 0.3] (p1)--(P1)--(P2)--(p2)--cycle;
	
	% --
	}
	
	% labels
	{ [canvas is yz plane at x = \A]
	
	\draw[gray!50]
		(-\A,-0.4) coordinate (A1) edge ++(0,0.1) edge ++(0,-0.1)
		--
		(\A,-0.4) coordinate (A2) edge ++(0,0.1) edge ++(0,-0.1);
	
	\path (A1)--node[inner sep = 2pt, fill = white] {$B$} (A2);
	
	}
	
	{ [canvas is yz plane at x = -\a]
	
	\draw[gray!50]
		(-\a,\H + 0.4) coordinate (a1) edge ++(0,0.1) edge ++(0,-0.1)
		--
		(\a,\H + 0.4) coordinate (a2) edge ++(0,0.1) edge ++(0,-0.1);
	
	\path (a1) -- node[inner sep = 2pt, fill = white] {$b$} (a2);
	
	}
	
	\draw[dashed] (O) --node[left] {$H$} (0,0,\H) coordinate (o);
	\fill
		(O) circle (1pt)
		(o) circle (1pt);
	
	% --
	}

	% --	
	\end{scope}
	
	\end{scope}

\end{tikzpicture}	
\end{document}